Because the early-warning and BN layers consume overlapping factor indicators, one might worry that their agreement is tautological. The experimental design reduces that risk in two ways. First, the champion is trained to predict the eventual severe label from prefix statistics, not to imitate BN posteriors; its strong TEST AUC therefore reflects schedule dynamics and covariates beyond binary factor flags. Second, BN policies are evaluated as inspection procedures under budget, not as classifiers trained on the same prefix matrix. When both layers are useful, they are useful for different decisions: triage versus explainable evidence collection~\cite{Egwim2021,Luu2009,Memarzadeh2016}.

The choice of HistGradientBoosting rather than a deep sequence model follows from the tabular nature of the released features and from literature showing competitive performance of boosted trees on structured project attributes~\cite{Egwim2021,Friedman2001,Tagharobi2025}. Sequence models could in principle exploit ordered update tokens, but they would require a different leakage analysis and are outside the executed campaign. Staying within tabular boosting keeps the method comparable to common delay-ML baselines and keeps training reproducible on commodity hardware.

Parameter provenance for $(p_c,\ell)$ and for the cost ratio deserves restatement because these values influence posterior levels. The cause probability $p_c{=}0.55$ and leak $\ell{=}0.02$ are the TRAIN-estimated configuration used throughout the BN replay artifacts; they are not borrowed from medical Noisy-OR libraries~\cite{Onisko2001,Heckerman1996}. The $5{:}1$ miss-to-false-alarm ratio is a design choice encoding higher miss cost for severe delay alerts, analogous in spirit to asymmetric clinical utilities but not claimed as an empirical industry constant~\cite{Vickers2006,Howard1966}. If an agency’s relative costs differ, both VoI rankings and net-benefit curves should be recomputed; the software stack is built for that recomputation.

Dataset scale should be read in context. With $3661$ PIDs and $550$ TEST projects, the study is larger than many elicitation-only BN delay case studies and comparable to recent ML delay papers that use hundreds to thousands of projects~\cite{Khodabakhshian2025,Egwim2021,Tagharobi2025}. It is not a multi-country mega-panel. Missingness in free-text reasons and uneven agency representation are real biases. The integrity response is documentation and PID-held-out evaluation, not a claim that NYC Open Data is a universal population~\cite{Woldesenbet2016,Flyvbjerg2009}.

Taken together, the executed evidence supports a restrained thesis: calibrated prefix-based early warning on open capital-project histories can exceed commonly cited accuracy references while remaining probabilistically usable, and budgeted Bayesian inspection provides an auditable factor-level account whose policy variants may tie when reason evidence is sparse. That thesis is narrower than a claim of universal VoI superiority, and it is stronger than a purely methodological proposal without TEST numbers~\cite{Song2021,Gneiting2007,Luu2009,Tagharobi2025}.
